% GENERATED FILE. Edit canonical lessons, then run npm run generate:books.
% canonical-source-hash: fnv1a64:de88c92de03b1f68
% canonical-lessons: HI-C487-chut, HI-C487-sankraman, HI-C487-khasra, HI-C487-cecak, HI-C487-pakna

\chapter{Catching Something}
\label{ch:hi-catching-something}
\hlchaptermodality{487}

\begin{chapteropening}
\textbf{By the end of this chapter:} \emph{I can talk about illnesses that spread.}

Complete the last lesson of chapter 487: I can talk about illnesses that spread.
\end{chapteropening}

\section[\texorpdfstring{\hi{छूत} (chūt)}{छूत (chūt)}]{\hi{छूत} (chūt) — contagion; infection by touch}

\label{lesson:HI-C487-chut}



\begin{quote}
Before the new one: say the Hindi for the spine, then the Hindi for to swell (of the body).
\end{quote}

\subsection*{You'll want to know: \hi{छूत}}
\textbf{\hi{छूत}} — \emph{chūt} — \textquotedblleft{}contagion; infection by touch\textquotedblright{}.

\begin{quote}
\hi{छूत} \hi{की} \hi{बीमारी} \hi{में} \hi{मरीज़} \hi{को} \hi{अलग} \hi{रखते} \hi{हैं।}
\end{quote}

\begin{quote}
\emph{chūt kī bīmārī meṁ marīz ko alag rakhte haiṁ.}
\end{quote}

\begin{quote}
With a contagious illness, the patient is kept apart.
\end{quote}

Feminine.

\begin{cousinweb}
\textbf{\hi{छूत}} (\emph{chūt}) is inherited: CDIAL lists Hindi \emph{chūt}, \textquotedblleft{}touch of anything impure\textquotedblright{}, under \emph{*chupti}, a form it reconstructs (CDIAL 5057), from the same root as \emph{chūnā}, \textquotedblleft{}to touch\textquotedblright{}. \emph{chūt kī bīmārī} is a disease that spreads by contact.
\end{cousinweb}

\begin{grammarlens}[title={What you've built}]
One more word for describing what you lived and felt.
\end{grammarlens}

\subsection*{Your turn}
\begin{itemize}
\raggedright
  \item \emph{Say it:} \emph{chūt}
  \item \emph{Say it:} \emph{chūt}, once more
  \item \emph{Say it:} \emph{chūt kī bīmārī meṁ marīz ko alag rakhte haiṁ.} — With a contagious illness, the patient is kept apart.
  \item \emph{Recall:} say the Hindi for the spine, then the Hindi for to swell (of the body), then say \emph{chūt} again
\end{itemize}

\begin{tcolorbox}[breakable,colback=teal!4,colframe=teal!35!black,title={Before you move on}]
What does \hi{छूत} mean? (\textquotedblleft{}Contagion; infection by touch\textquotedblright{}.) Say it once more.
\end{tcolorbox}
\section[\texorpdfstring{\hi{संक्रमण} (saṅkramaṇ)}{संक्रमण (saṅkramaṇ)}]{\hi{संक्रमण} (saṅkramaṇ) — an infection}

\label{lesson:HI-C487-sankraman}



\begin{quote}
Before the new one: say the Hindi for to swell (of the body), then the Hindi for contagion.
\end{quote}

\subsection*{You'll want to know: \hi{संक्रमण}}
\textbf{\hi{संक्रमण}} — \emph{saṅkramaṇ} — \textquotedblleft{}an infection\textquotedblright{}.

\begin{quote}
\hi{गंदे} \hi{घाव} \hi{में} \hi{संक्रमण} \hi{हो} \hi{सकता} \hi{है।}
\end{quote}

\begin{quote}
\emph{gande ghāv meṁ saṅkramaṇ ho saktā hai.}
\end{quote}

\begin{quote}
A dirty wound can become infected.
\end{quote}

Masculine.

\begin{cousinweb}
\textbf{\hi{संक्रमण}} (\emph{saṅkramaṇ}) is a tatsama, borrowed from Sanskrit \emph{saṁkramaṇa}, \textquotedblleft{}going together, entrance into, transference\textquotedblright{} (Monier-Williams). The medical sense, a disease passing into the body, is modern.
\end{cousinweb}

\begin{grammarlens}[title={What you've built}]
One more word for describing what you lived and felt.
\end{grammarlens}

\subsection*{Your turn}
\begin{itemize}
\raggedright
  \item \emph{Say it:} \emph{saṅkramaṇ}
  \item \emph{Say it:} \emph{saṅkramaṇ}, once more
  \item \emph{Say it:} \emph{gande ghāv meṁ saṅkramaṇ ho saktā hai.} — A dirty wound can become infected.
  \item \emph{Recall:} say the Hindi for to swell (of the body), then the Hindi for contagion, then say \emph{saṅkramaṇ} again
\end{itemize}

\begin{tcolorbox}[breakable,colback=teal!4,colframe=teal!35!black,title={Before you move on}]
What does \hi{संक्रमण} mean? (\textquotedblleft{}An infection\textquotedblright{}.) Say it once more.
\end{tcolorbox}
\section[\texorpdfstring{\hi{खसरा} (khasrā)}{खसरा (khasrā)}]{\hi{खसरा} (khasrā) — measles}

\label{lesson:HI-C487-khasra}



\begin{quote}
Before the new one: say the Hindi for contagion, then the Hindi for an infection.
\end{quote}

\subsection*{You'll want to know: \hi{खसरा}}
\textbf{\hi{खसरा}} — \emph{khasrā} — \textquotedblleft{}measles\textquotedblright{}.

\begin{quote}
\hi{मेरे} \hi{भाई} \hi{को} \hi{बचपन} \hi{में} \hi{खसरा} \hi{हुआ} \hi{था।}
\end{quote}

\begin{quote}
\emph{mere bhāī ko bacpan meṁ khasrā huā thā.}
\end{quote}

\begin{quote}
My brother had measles as a child.
\end{quote}

Masculine.

\begin{cousinweb}
\textbf{\hi{खसरा}} (\emph{khasrā}) is a tadbhava: CDIAL lists Hindi \emph{khasrā}, \textquotedblleft{}scab, measles\textquotedblright{}, under Sanskrit \emph{khasa} (CDIAL 3854).
\end{cousinweb}

\begin{grammarlens}[title={What you've built}]
One more word for describing what you lived and felt.
\end{grammarlens}

\subsection*{Your turn}
\begin{itemize}
\raggedright
  \item \emph{Say it:} \emph{khasrā}
  \item \emph{Say it:} \emph{khasrā}, once more
  \item \emph{Say it:} \emph{mere bhāī ko bacpan meṁ khasrā huā thā.} — My brother had measles as a child.
  \item \emph{Recall:} say the Hindi for contagion, then the Hindi for an infection, then say \emph{khasrā} again
\end{itemize}

\begin{tcolorbox}[breakable,colback=teal!4,colframe=teal!35!black,title={Before you move on}]
What does \hi{खसरा} mean? (\textquotedblleft{}Measles\textquotedblright{}.) Say it once more.
\end{tcolorbox}
\section[\texorpdfstring{\hi{चेचक} (cecak)}{चेचक (cecak)}]{\hi{चेचक} (cecak) — smallpox}

\label{lesson:HI-C487-cecak}



\begin{quote}
Before the new one: say the Hindi for an infection, then the Hindi for measles.
\end{quote}

\subsection*{You'll want to know: \hi{चेचक}}
\textbf{\hi{चेचक}} — \emph{cecak} — \textquotedblleft{}smallpox\textquotedblright{}.

\begin{quote}
\hi{टीके} \hi{की} \hi{वजह} \hi{से} \hi{अब} \hi{चेचक} \hi{नहीं} \hi{होती।}
\end{quote}

\begin{quote}
\emph{ṭīke kī vajah se ab cecak nahīṃ hotī.}
\end{quote}

\begin{quote}
Thanks to vaccination, smallpox no longer occurs.
\end{quote}

Feminine.

\begin{cousinweb}
\textbf{\hi{चेचक}} (\emph{cecak}) is usually traced to Persian \emph{chēchak}, \textquotedblleft{}smallpox\textquotedblright{}.
\end{cousinweb}

\begin{grammarlens}[title={What you've built}]
One more word for describing what you lived and felt.
\end{grammarlens}

\subsection*{Your turn}
\begin{itemize}
\raggedright
  \item \emph{Say it:} \emph{cecak}
  \item \emph{Say it:} \emph{cecak}, once more
  \item \emph{Say it:} \emph{ṭīke kī vajah se ab cecak nahīṃ hotī.} — Thanks to vaccination, smallpox no longer occurs.
  \item \emph{Recall:} say the Hindi for an infection, then the Hindi for measles, then say \emph{cecak} again
\end{itemize}

\begin{tcolorbox}[breakable,colback=teal!4,colframe=teal!35!black,title={Before you move on}]
What does \hi{चेचक} mean? (\textquotedblleft{}Smallpox\textquotedblright{}.) Say it once more.
\end{tcolorbox}
\section[\texorpdfstring{\hi{पकना} (paknā)}{पकना (paknā)}]{\hi{पकना} (paknā) — to fester; to ripen}

\label{lesson:HI-C487-pakna}



\begin{quote}
Before the new one: say the Hindi for measles, then the Hindi for smallpox.
\end{quote}

\subsection*{You'll want to know: \hi{पकना}}
\textbf{\hi{पकना}} — \emph{paknā} — \textquotedblleft{}to fester; to ripen\textquotedblright{}.

\begin{quote}
\hi{पुराना} \hi{घाव} \hi{पक} \hi{गया} \hi{और} \hi{बहुत} \hi{दर्द} \hi{हुआ।}
\end{quote}

\begin{quote}
\emph{purānā ghāv pak gayā aur bahut dard huā.}
\end{quote}

\begin{quote}
The old wound festered and hurt a lot.
\end{quote}

Past (it …): \hi{वह} \hi{घाव} \hi{पक} \hi{गया}, \emph{vah ghāv pak gayā}.

\begin{cousinweb}
\textbf{\hi{पकना}} (\emph{paknā}) is a tadbhava: CDIAL lists Hindi \emph{paknā}, \textquotedblleft{}to be cooked, become ripe, suppurate\textquotedblright{}, under Sanskrit \emph{pakvá}, \textquotedblleft{}cooked, ripe\textquotedblright{} (CDIAL 7621). Fruit ripens and a wound festers with the same verb; \emph{pakānā}, \textquotedblleft{}to cook\textquotedblright{}, is its causative.
\end{cousinweb}

\begin{grammarlens}[title={What you've built}]
One more word for describing what you lived and felt.
\end{grammarlens}

\subsection*{Your turn}
\begin{itemize}
\raggedright
  \item \emph{Say it:} \emph{paknā}
  \item \emph{Say it:} \emph{paknā}, once more
  \item \emph{Say it:} \emph{purānā ghāv pak gayā aur bahut dard huā.} — The old wound festered and hurt a lot.
  \item \emph{Recall:} say the Hindi for measles, then the Hindi for smallpox, then say \emph{paknā} again
\end{itemize}

\begin{tcolorbox}[breakable,colback=teal!4,colframe=teal!35!black,title={Before you move on}]
What does \hi{पकना} mean? (\textquotedblleft{}To fester; to ripen\textquotedblright{}.) Say it once more.
\end{tcolorbox}
